% Figura: balanza de torsión de Coulomb (esquema).
\begin{figure}[ht]
\centering
\tikzset{esf/.style={circle, inner sep=0pt, minimum size=9pt, ball color=azulacero!75, draw=azulprof, line width=0.4pt},
         esffija/.style={circle, inner sep=0pt, minimum size=9pt, ball color=acento!80, draw=acento!60!black, line width=0.4pt}}
\begin{minipage}[t]{0.56\textwidth}\centering
\begin{tikzpicture}[baseline=(current bounding box.north), x=1cm, y=1cm]
  % (a) vista en perspectiva: cilindro de vidrio, hilo, aguja y esferas
  \def\rx{1.55}\def\ry{0.42}
  \fill[azulclaro!30] (-\rx,0) -- (-\rx,3.2) arc (180:360:{\rx} and \ry) -- (\rx,0) arc (0:-180:{\rx} and \ry) -- cycle;
  \draw[azulprof, line width=0.8pt] (-\rx,0) -- (-\rx,3.2) (\rx,0) -- (\rx,3.2);
  \draw[azulprof, line width=0.8pt, fill=azulclaro!45] (0,3.2) ellipse ({\rx} and \ry);
  \draw[azulprof, line width=0.8pt] (-\rx,0) arc (180:360:{\rx} and \ry);
  \draw[azulprof, line width=0.5pt, densely dashed] (-\rx,0) arc (180:0:{\rx} and \ry);
  % escala graduada alrededor del cilindro (marcas en la mitad frontal)
  \draw[azulprof!70, line width=0.5pt] (-\rx,1.1) arc (180:360:{\rx} and \ry);
  \foreach \a in {190,200,...,350}{\draw[azulprof!70, line width=0.35pt] ({\rx*cos(\a)},{1.1+\ry*sin(\a)}) -- ({\rx*cos(\a)},{1.1+\ry*sin(\a)+0.1});}
  % cabezal micrométrico
  \draw[azulprof, line width=0.7pt, fill=azulclaro!70] (-0.45,3.52) rectangle (0.45,3.95);
  \draw[azulprof, line width=0.7pt, fill=azulclaro!40] (0,3.95) ellipse (0.45 and 0.13);
  \draw[acento, line width=0.9pt] (0,3.95) -- (0.3,3.99);
  \draw[azulprof, line width=0.7pt, fill=azulclaro!55] (0,3.52) ellipse (0.45 and 0.13);
  % hilo de torsión
  \draw[grispizarra, line width=0.6pt] (0,3.38) -- (0,1.45);
  % aguja (en perspectiva) con esfera móvil y contrapeso
  \draw[grispizarra, line width=1.3pt] (-1.15,1.28) -- (1.15,1.62);
  \fill[azulprof] (0,1.45) circle (1.3pt);
  \node[esf] at (1.15,1.62) {}; \node[esf] at (-1.15,1.28) {};
  % esfera fija sujeta desde la tapa
  \draw[grispizarra, line width=0.9pt] (1.02,3.05) -- (1.02,1.35);
  \node[esffija, minimum size=8pt] at (1.02,1.30) {};
  % rótulos
  \draw[grispizarra, line width=0.3pt] (0.5,3.8) -- (2.3,3.8); \node[font=\footnotesize, anchor=west] at (2.35,3.8) {cabezal micrométrico};
  \draw[grispizarra, line width=0.3pt] (0.02,2.6) -- (2.3,2.6); \node[font=\footnotesize, anchor=west] at (2.35,2.6) {hilo de torsión};
  \draw[grispizarra, line width=0.3pt] (1.25,1.7) -- (2.3,1.95); \node[font=\footnotesize, anchor=west] at (2.35,1.95) {esfera móvil};
  \draw[grispizarra, line width=0.3pt] (1.12,1.18) -- (2.3,1.05); \node[font=\footnotesize, anchor=west] at (2.35,1.05) {esfera fija};
  \draw[grispizarra, line width=0.3pt] (1.5,0.62) -- (2.3,0.35); \node[font=\footnotesize, anchor=west] at (2.35,0.3) {escala graduada};
  \node[font=\footnotesize, anchor=east] at (-1.75,1.1) {aguja};
  \draw[grispizarra, line width=0.3pt] (-1.7,1.15) -- (-1.25,1.27);
\end{tikzpicture}\\[2pt]{\small\scshape (a)}
\end{minipage}\hfill
\begin{minipage}[t]{0.40\textwidth}\centering
\begin{tikzpicture}[baseline=(current bounding box.north), x=1cm, y=1cm]
  % (b) vista superior: escala graduada y ángulo de torsión
  \def\Rg{1.65}
  \draw[azulprof, line width=0.8pt] (0,0) circle (\Rg);
  \foreach \a in {0,6,...,354}{\draw[azulprof!70, line width=0.3pt] (\a:\Rg) -- (\a:{\Rg-0.09});}
  \foreach \a in {0,30,...,330}{\draw[azulprof, line width=0.5pt] (\a:\Rg) -- (\a:{\Rg-0.17});}
  \draw[verdeteo, densely dashed, line width=0.5pt] (0,0) -- (0:1.3);
  \draw[grispizarra, line width=1.3pt] (200:1.2) -- (22:1.2);
  \fill[azulprof] (0,0) circle (1.4pt) node[below left, font=\scriptsize] {$O$};
  \node[esf] at (22:1.2) {}; \node[esf] at (200:1.2) {};
  \node[esffija] at (0:1.3) {};
  \draw[verdeteo, line width=0.6pt] (0:0.7) arc (0:22:0.7);
  \node[font=\small, text=verdeteo] at (11:0.95) {$\theta$};
  \node[font=\scriptsize, text=acento, anchor=west] at (1.85,-0.35) {fija};
  \node[font=\scriptsize, text=azulprof, anchor=south west] at (1.35,1.25) {móvil};
\end{tikzpicture}\\[2pt]{\small\scshape (b)}
\end{minipage}
\caption{Balanza de torsión de Coulomb (esquema). (a)~Una aguja aislante suspendida de un hilo fino lleva una esfera móvil en un extremo y un contrapeso en el otro; una segunda esfera
fija se introduce desde la tapa. Un cabezal micrométrico permite girar el extremo superior del hilo, y una escala graduada rodea el cilindro de vidrio. (b)~Vista desde arriba: al cargar ambas esferas, la repulsión desvía la aguja un
ángulo $\theta$ medido sobre la escala; en el equilibrio, el par recuperador del hilo, proporcional a $\theta$, iguala al de la fuerza electrostática. Girando el cabezal se varía la distancia entre las esferas y se comprueba cómo escala la fuerza con ella.}
\label{fig:torsion}
\end{figure}
